\chapter{Borg: An Open-Source TBDR GPU}
\label{chap:borg}

\section{Design Goals}

\borg{} is designed around three principles~\cite{Wendleder2024Borg}:
\begin{enumerate}
  \item \textbf{Open.} Every component -- HDL source, toolchain, tapeout scripts -- is publicly available.
  \item \textbf{Verifiable.} The implementation runs on both FPGA and real silicon.
  \item \textbf{Pedagogical.} The architecture is legible: a student can read and understand every module.
\end{enumerate}

\section{Tile-Based Deferred Rendering Architecture}

\borg{} follows the TBDR paradigm pioneered by PowerVR~\cite{Img1996TBDR}
and later adopted by Apple~GPU, ARM~Mali, and Vivante.
Geometry is processed tile-by-tile; only visible fragments are shaded,
reducing external memory bandwidth significantly.

\subsection{Morton-Encoded Tile Buffer}

The tile buffer uses Morton (Z-curve) encoding for cache-friendly 2D access:
\begin{equation}
  \text{addr}(x, y) = \text{interleave\_bits}(x, y)
\end{equation}
\todo{Explain interleaving, cache line utilisation, comparison to linear layout.}

\subsection{Rasterizer}

\todo{Edge function rasteriser, subpixel precision, tile dispatch.}

\subsection{Shader Execution Unit}

\todo{Custom ISA, register file, ALU, branch handling.}

\section{Custom ISA}

\todo{Instruction set reference table. Encoding, decode pipeline.}

\section{FPGA Implementation}

\subsection{iCE40 Target}

\borg{} is synthesized for the Lattice iCE40 family using the open-source
toolchain (Yosys + nextpnr).
The design fits within the iCE40~UP5K (5280~logic cells).

\subsection{Resource Utilisation}

\begin{table}[h]
  \centering
  \begin{tabular}{lrr}
    \toprule
    Resource & Used & Available \\
    \midrule
    Logic Cells & \todo{N} & 5280 \\
    BRAMs       & \todo{N} & 30 \\
    \bottomrule
  \end{tabular}
  \caption{iCE40 UP5K resource utilisation.}
  \label{tab:fpga_util}
\end{table}

\section{ASIC Tapeout}

\borg{} has been taped out via the Tiny Tapeout programme in two process nodes:

\begin{description}
  \item[TTIHP26a (IHP SG13G2, 130\,nm)] Floating-point adder peripheral.
    March 2026. Chips expected July 2026.
  \item[TTSKY26a (SkyWater SKY130, 130\,nm)] Triangle rasteriser with
    vertex shader. May 2026.
  \item[GF180 via Efabless] Diffuse shader ASIC.
    Funded by the Google Open-Source Silicon Programme. 2022.
\end{description}

\section{Toolchain}

\begin{description}
  \item[HDL] Chisel (Scala-based), compiled via FIRRTL to Verilog.
  \item[Synthesis] Yosys.
  \item[Place \& Route] nextpnr (FPGA), OpenROAD (ASIC).
  \item[Simulation] Verilator, cocotb.
  \item[Register Map] PeakRDL-chisel.
\end{description}

\section{NLnet Funding}

The \borg{} project is funded by the
NGI0 Commons Fund, a fund established by NLnet Foundation with financial
support from the European Commission's Next Generation Internet programme
(Grant Agreement No.\ 101135429, DG CONNECT), and by the Swiss State
Secretariat for Education, Research and Innovation (SERI).
